% ============================================================
%  PLANTILLA DE MAPA MENTAL — syntheca
%  Clase STANDALONE/ARTICLE, usa la librería `mindmap` de TikZ
%  (distribución radial automática, no hay que calcular ángulos
%  a mano por nodo).
%
%  Filosofía de color, distinta al resto del sistema: acá SÍ se
%  usa color para diferenciar ramas (familias de conceptos) —
%  no es decoración, es una ayuda de codificación dual (dual
%  coding) para la memoria, la misma lógica por la que
%  pedagogia-cognitiva permite diagramas cuando cumplen una
%  función real. La paleta es apagada/profesional, nunca
%  saturada — mantiene el mismo espíritu sobrio del resto.
%
%  REGENERACIÓN COMPLETA: a diferencia del libro (que acumula
%  capítulos), este archivo se reescribe ENTERO cada vez que
%  generador-mapa-mental corre — nunca se edita a mano un nodo
%  suelto. La fuente de verdad es mapa-estudio.json, no el .tex.
% ============================================================
\usepackage[utf8]{inputenc}
\usepackage[T1]{fontenc}
\usepackage{tikz}
\usetikzlibrary{mindmap,trees,backgrounds}
\usepackage{fancyhdr}
\usepackage[table]{xcolor}
\usepackage{geometry}
\geometry{top=1.2cm, bottom=1.2cm, left=1cm, right=1cm, landscape}

\definecolor{textgray}{RGB}{60, 60, 60}
\definecolor{rulegray}{RGB}{150, 150, 150}

% ---------- Paleta de ramas — apagada, profesional, máximo 6 colores.
% Si un examen tiene más de 6 familias, se reciclan en orden. ----------
\definecolor{ramaUno}{RGB}{91, 110, 130}    % azul-gris apagado
\definecolor{ramaDos}{RGB}{120, 100, 90}    % marrón-gris apagado
\definecolor{ramaTres}{RGB}{95, 120, 100}   % verde-gris apagado
\definecolor{ramaCuatro}{RGB}{115, 95, 120} % violeta-gris apagado
\definecolor{ramaCinco}{RGB}{130, 110, 80}  % ocre apagado
\definecolor{ramaSeis}{RGB}{90, 115, 120}   % teal-gris apagado

\newcommand{\miNombre}{Nombre Apellido}
\newcommand{\materiaActual}{Materia}
\newcommand{\setMateria}[1]{\renewcommand{\materiaActual}{#1}}
\newcommand{\nombreMapa}{Mapa mental}
\newcommand{\setNombreMapa}[1]{\renewcommand{\nombreMapa}{#1}}

\pagestyle{fancy}
\fancyhf{}
\renewcommand{\headrulewidth}{0pt}
\renewcommand{\footrulewidth}{0pt}
\fancyfoot[C]{\small\color{textgray}\thepage}
\fancyfoot[R]{\footnotesize\color{textgray}\miNombre}
\fancyfoot[L]{\footnotesize\color{textgray}\materiaActual}

% ---------- Estilo base de los conceptos (nodos) ----------
% Nivel raíz: el examen/tema central.
% Nivel 1: familias de conceptos (ramas, coloreadas).
% Nivel 2: conceptos individuales (hojas, color heredado de su rama).
\tikzstyle{every node}=[font=\sffamily]
\tikzset{
  root concept/.style={
    concept color=textgray,
    font=\sffamily\bfseries\Large,
    text=white,
    minimum size=3.2cm,
  },
  level 1 concept/.append style={
    font=\sffamily\bfseries\small,
    text=white,
    sibling angle=90,
    level distance=4.6cm,
    minimum size=2.1cm,
  },
  level 2 concept/.append style={
    font=\sffamily\scriptsize,
    text=white,
    sibling angle=45,
    level distance=2.6cm,
    minimum size=1.5cm,
  },
}
